
\documentclass{article}
\usepackage{colortbl}
\usepackage{makecell}
\usepackage{multirow}
\usepackage{supertabular}

\begin{document}

\newcounter{utterance}

\centering \large Interaction Transcript for game `eqbench', experiment `eq\_bench', episode 85 with reflexion{-}Qwen3.5{-}27B.
\vspace{24pt}

{ \footnotesize  \setcounter{utterance}{1}
\setlength{\tabcolsep}{0pt}
\begin{supertabular}{c@{$\;$}|p{.15\linewidth}@{}p{.15\linewidth}p{.15\linewidth}p{.15\linewidth}p{.15\linewidth}p{.15\linewidth}}
    \# & \multicolumn{2}{c}{Player} && \multicolumn{2}{c}{Game Master} \\
    \hline

    \theutterance \stepcounter{utterance}  
    & & & \multicolumn{4}{p{0.6\linewidth}}{
        \cellcolor[rgb]{0.9,0.9,0.9}{
            \makecell[{{p{\linewidth}}}]{
                \texttt{\tiny{[P1$\langle$GM]}}
                \\ 
\texttt{Your task is to predict the likely emotional responses of a character in this dialogue:} \\
\\ 
\texttt{Charlie: You're too green for this, Ruiz.} \\
\texttt{Mari: And you're too washed up, Redburn.} \\
\texttt{Charlie: I've been on this beat longer than you've been alive, kid.} \\
\texttt{Mari: And yet, here we are, stuck on the same case. Seems like experience isn't everything.} \\
\texttt{Charlie: Don't get smart with me, Mari.} \\
\texttt{Mari: Someone has to be the smart one here.} \\
\texttt{Charlie: You think I don't know what you're doing? Trying to prove yourself, make a name?} \\
\texttt{Mari: Better than drinking myself into oblivion.} \\
\texttt{Charlie: That's a low blow, Ruiz.} \\
\texttt{Mari: Truth hurts, doesn't it?} \\
\texttt{Charlie: You don't know a damn thing about me.} \\
\texttt{Mari: And you don't know anything about me, either. So maybe we should just focus on the case.} \\
\texttt{Charlie: Maybe you should focus on following protocol for once.} \\
\texttt{Mari: And maybe you should try embracing the 21st century.} \\
\texttt{Charlie: I don't need some fancy tech to do my job.} \\
\texttt{Mari: And I don't need some old man telling me how to do mine.} \\
\texttt{Charlie: You're out of line, Ruiz.} \\
\texttt{Mari: And you're out of touch, Redburn.} \\
\texttt{{[}End dialogue{]}} \\
\\ 
\texttt{At the end of this dialogue, Mari would feel...} \\
\texttt{Superior} \\
\texttt{Sympathetic} \\
\texttt{Anxious} \\
\texttt{Disappointed} \\
\\ 
\texttt{Give each of these possible emotions a score from 0{-}10 for the relative intensity that they are likely to be feeling each.} \\
\\ 
\texttt{You must output in the following format, including headings (of course, you should give your own scores), with no additional commentary:} \\
\\ 
\texttt{Superior: <score>} \\
\texttt{Sympathetic: <score>} \\
\texttt{Anxious: <score>} \\
\texttt{Disappointed: <score>} \\
\\ 
\\ 
\texttt{{[}End of answer{]}} \\
\\ 
\texttt{Remember: zero is a valid score, meaning they are likely not feeling that emotion. You must score at least one emotion > 0.} \\
\\ 
\texttt{Your answer:} \\
            }
        }
    }
     \\ \\

    \theutterance \stepcounter{utterance}  
    & \multicolumn{4}{p{0.6\linewidth}}{
        \cellcolor[rgb]{0.561,0.737,0.561}{
            \makecell[{{p{\linewidth}}}]{
                \texttt{\tiny{[P1$\rangle$GM]}}
                \texttt{Superior: 8} \\
\texttt{Sympathetic: 0} \\
\texttt{Anxious: 1} \\
\texttt{Disappointed: 2} \\
            }
        }
    }
    & & \\ \\

    \theutterance \stepcounter{utterance}  
    & & & \multicolumn{2}{p{0.3\linewidth}}{
        \cellcolor[rgb]{0.9,0.9,0.9}{
            \makecell[{{p{\linewidth}}}]{
                \texttt{\tiny{[GM$|$GM]}}
                \texttt{Superior: 8} \\
\texttt{Sympathetic: 0} \\
\texttt{Anxious: 1} \\
\texttt{Disappointed: 2} \\
            }
        }
    }
    & & \\ \\

    \theutterance \stepcounter{utterance}  
    & & & \multicolumn{2}{p{0.3\linewidth}}{
        \cellcolor[rgb]{0.9,0.9,0.9}{
            \makecell[{{p{\linewidth}}}]{
                \texttt{\tiny{[GM$|$GM]}}
                \texttt{Superior: 8} \\
\texttt{Sympathetic: 0} \\
\texttt{Anxious: 0} \\
\texttt{Disappointed: 4} \\
            }
        }
    }
    & & \\ \\

    \theutterance \stepcounter{utterance}  
    & & & \multicolumn{2}{p{0.3\linewidth}}{
        \cellcolor[rgb]{0.9,0.9,0.9}{
            \makecell[{{p{\linewidth}}}]{
                \texttt{\tiny{[GM$|$GM]}}
                \texttt{game\_result = LOSE} \\
            }
        }
    }
    & & \\ \\

    \theutterance \stepcounter{utterance}  
    & & \multicolumn{4}{p{0.6\linewidth}}{
        \cellcolor[rgb]{0.95,0.95,0.95}{
            \makecell[{{p{\linewidth}}}]{
                \texttt{\tiny{[P1$\leftrightarrow$P]}}
                \texttt{Post{-}episode Reflexion (not sent to GM):} \\
\texttt{New plan: My previous attempt failed because I assigned a high score to "Superior" (8) without sufficiently weighing the context of the final exchange. While Mari is confident, her closing line ("And you're out of touch, Redburn") is a direct rebuttal to Charlie's accusation of being "out of line," suggesting her primary emotional state is defensive and frustrated rather than feeling a sense of superiority. Additionally, I underestimated "Disappointed" (2); the dialogue reveals a deep professional rift where Mari feels let down by Charlie's refusal to adapt ("embracing the 21st century") and his reliance on outdated methods. Next time, I will analyze the final lines specifically to determine if the character is asserting dominance (Superior) or expressing frustration with the other's limitations (Disappointed). I will likely lower the "Superior" score and significantly raise the "Disappointed" score to reflect the futility of the argument and the breakdown in their working relationship.} \\
            }
        }
    }
    & \\ \\

\end{supertabular}
}

\end{document}
